\chapter*{Abstract}

Located in the Pampa Amarilla in Argentina, the Pierre Auger Observatory combines different detection techniques to study ultra-high-energy cosmic rays in a multi-hybrid approach.
To extend the measurements towards lower energies, two denser surface-detector arrays, the \sd{750} and \sd{433}, were deployed within the 1500\,m array, with spacings of 750\,m and 433\,m, respectively.
This work focuses on the \sd{433}, which allows the cosmic-ray energy spectrum to be measured from $63\,\mathrm{PeV}$ using water-Cherenkov detectors.
Its energy range allows the characterization of the steepening of the spectrum known as the second knee.
A complete characterization of this feature is key to understanding the transition from Galactic to extragalactic cosmic rays.

The analysis uses \SDFourThreeThreeNEvents{} events recorded between July 2018 and December 2021, with zenith angles below $40^\circ$.
Above the selected energy threshold, the array can be treated as fully efficient, allowing the exposure to be calculated geometrically without relying on simulations.

Given the limited number of events measured simultaneously by the \sd{433} and the fluorescence detectors, a direct energy calibration is not possible.
However, events measured by both the \sd{433} and \sd{750} allow the \sd{433} shower size to be calibrated against the \sd{750} energy.
In this way, the energy scale of the fluorescence detector is inherited by the \sd{433}.
With the primary energy reconstructed for each event, the raw spectrum is then derived.
To measure the second knee, a likelihood fit is performed folding the spectrum model with the detector response.

The spectrum is described by a broken power law with a smooth transition corresponding to the second knee.
Using \sd{433} data only, the second knee is measured at $E_{01}=(2.8\pm0.5_{\mathrm{stat}}\pm0.7_{\mathrm{syst}})\times10^{17}\,\mathrm{eV}$ and the spectral index changes from $\gamma_0=3.04\pm0.02_{\mathrm{stat}}\pm0.09_{\mathrm{syst}}$ below the break to $\gamma_1=3.34\pm0.08_{\mathrm{stat}}\pm0.11_{\mathrm{syst}}$ above it.
The fluorescence-detector energy scale dominates the systematic uncertainty in the flux and the feature position, while the calibration uncertainty has the largest effect on the spectral indices.

An update of the \sd{750} spectrum published in 2021 is also presented using the current energy scale of the Observatory.
With the events recorded by both arrays, the \sd{750} detection efficiency is also re-estimated.
Since the second-knee position cannot be measured using \sd{750} data only, a preliminary combined spectrum is obtained using data from both arrays.
With this combination, the aim is to better constrain the second-knee position and extend the spectrum measurement through the ankle, thus helping to understand the transition from Galactic to extragalactic cosmic rays.
